
This study set out to test whether stable spectral properties indicate model quality---and obtained the opposite result. The coefficient of variation of participation ratio (CV\_PR), a measure of randomized SVD estimator variance, was hypothesized to correlate negatively with ImageNet accuracy: models with flat spectra should produce consistent singular value estimates across random projections, and this stability should signal better generalization. Experiments on 100 pretrained models from the timm library revealed a strong positive correlation (r = +0.61, p < $10^{-10})$.

Spectral analysis of neural network weights has become a method for assessing model quality without requiring evaluation data. Martin and Mahoney (2019) established that heavy-tailed eigenvalue distributions reflect training quality. Unterthiner et al. (2020) demonstrated that simple weight statistics predict accuracy with $R^2$ > 0.98. However, the directional relationship between spectral variance and generalization had not been systematically tested. The timm library hosts over 1,000 architectures with pretrained weights, and practitioners selecting models rely on implicit assumptions about spectral properties.

Prior work computed spectral features directly (effective rank, participation ratio, power-law exponents) but did not examine the variance of these estimates under randomized computation. Randomized SVD introduces stochasticity: projecting a weight matrix onto random subspaces yields different singular value approximations across seeds. The reasoning was that this variance reflects spectral shape---flat spectra should produce stable estimates, peaked spectra should not. No prior study had tested whether this computational variance predicts model quality.

The main finding is that CV\_PR positively correlates with accuracy. The assumption that stability signals quality does not hold for this metric. Two competing explanations are considered: (1) model size confounds both accuracy and spectral diversity---larger models have more parameters, higher accuracy, and potentially more varied spectral structure; (2) higher CV\_PR may reflect diverse feature representations rather than instability. Distinguishing these requires partial correlation analysis controlling for parameter count.

This work makes three contributions:

\begin{enumerate}
\item \textbf{Methodology validation}: CV\_PR can be reliably extracted from 100 pretrained models using randomized SVD with 20 seeds (100\% completion rate, CV\_PR values in [0.0014, 0.0326]).
\end{enumerate}

\begin{enumerate}
\item \textbf{Hypothesis falsification}: CV\_PR correlates positively with ImageNet accuracy (r = +0.6065, 95\% CI [0.506, 0.703]), contradicting the hypothesized negative correlation.
\end{enumerate}

\begin{enumerate}
\item \textbf{Interpretive framework}: Alternative explanations (confounding, inverted mechanism) are identified to guide future hypothesis reformulation.
\end{enumerate}
